% 数理物理1 数理物理1-10 問題7（072-parseval。当時は 数理物理1-8 問題6・shou08-l2-parseval） 解いた日 2026-10-03
%% kotae
\toi{1}{×}{未着手}
\mikai
\toi{2}{×}{未着手}
\mikai
\toi{3}{◎}{}
\begin{align*}
  &\frac{2}{L}\int_{-L/2}^{L/2}f(x)^{2}\dd x\\
  &\qquad=\frac{a_0^{2}}{2}+\sum_{1}^{\infty}\kakimachigai{a_n^{2}+b_n^{2}}{(a_n^{2}+b_n^{2})}
\end{align*}
\begin{align*}
  \frac{2d}{L}&=\frac{1}{2}\Bigl(\frac{2d}{L}\Bigr)^{2}\\
  &\qquad+\sum_{1}^{\infty}\Bigl(\frac{2\sin(\pi nd/L)}{\pi n}\Bigr)^{2}
\end{align*}
\begin{align*}
  \sum_{1}^{\infty}\frac{4}{\pi^{2}n^{2}}\sin^{2}\frac{\pi nd}{L}&=\frac{2d}{L}-\frac{2d^{2}}{L^{2}}\\
  \sum\frac{1}{n^{2}}\sin^{2}\frac{\pi nd}{L}&=\pi^{2}\Bigl(\frac{dL}{2L^{2}}-\frac{d^{2}}{2L^{2}}\Bigr)\\
  &=\pi^{2}\frac{dL-d^{2}}{2L^{2}}\\
  &=\frac{\pi^{2}d(L-d)}{2L^{2}}
\end{align*}
\toi{4}{◎}{}
\begin{align*}
  \sum\frac{1}{n^{2}}\sin^{2}\frac{\pi n}{2}&=\frac{\pi^{2}(L/2)^{2}}{2L^{2}}\\
  \sum\frac{1}{(2m+1)^{2}}&=\frac{\pi^{2}}{8}
\end{align*}
\begin{align*}
  \sum\frac{1}{n^{2}}\sin^{2}\frac{\pi n}{3}&=\frac{\pi^{2}\cdot\frac{L}{3}\cdot\frac{2L}{3}}{2L^{2}}=\frac{\pi^{2}}{9}\\
  \frac{3}{4}\sum\frac{1}{n^{2}}&=\frac{\pi^{2}}{9}
\end{align*}
\[ \frac{4\pi^{2}}{27} \]
%% kizuki
\item (1) は $\abs{f}^{2}=f^{*}f$ に展開を2つ代入して二重和にし，正規直交性 $\int_c^{c+L}\ee^{-\ii q_{\mu}x}\ee^{\ii q_{\nu}x}\dd x=L\delta_{\mu\nu}$ で一重の和にする。導出を問われた式は，使う前にこの1段を書く。
\item (2) は (1) の右辺を $\mu=0$ と $\mu=\pm n$ に分け，実関数では $\hat{f}_0=a_0/2$，$\abs{\hat{f}_{\pm n}}^{2}=(a_n^{2}+b_n^{2})/4$ を使うと出る（問題1（067）の関係）。
